\documentclass{itmo}

\setlist[itemize]{label=---}
\setlist[enumerate]{label=\arabic*)}

\usepackage{dcolumn}
\newcolumntype{d}[1]{D{.}{.}{#1}}

\begin{document}

\makereporttitle{}
	{лабораторной работе №5}
	{Структуры и алгоритмы в базах данных и распределенных системах}
	{Реализация обратного индекса}
	{}
	{Постнов~С.~А./P4135}
	{Платонов~А.~В./доцент, Портнов~П.~В./ассистент}

\maketableofcontents

\chapter{Теоретическая часть}

\section{Обратный индекс}

Обратный индекс (inverted index)~---~структура данных, отображающая термы в список документов (posting list), в которых они встречаются.
Для каждого терма хранится упорядоченный по \texttt{DocID} список постингов, содержащих идентификатор документа и позиции вхождения терма.
Обратный индекс является основной структурой данных в поисковых системах, обеспечивая быстрый поиск документов по ключевым словам.

Формат постинга: $(DocID, [pos_1, pos_2, \ldots, pos_k])$, где $DocID$~---~идентификатор документа, $pos_i$~---~позиции вхождения терма в документе.
Позиции необходимы для реализации операторов (\texttt{ADJ}, \texttt{NEAR}).

\section{Операции над posting lists}

Основные операции над posting lists, используемые в поисковых запросах:
\begin{itemize}
	\item \textit{Intersect} ($A \cap B$)~---~пересечение документов (оператор \texttt{AND});
	\item \textit{Union} ($A \cup B$)~---~объединение документов (оператор \texttt{OR});
	\item \textit{Difference} ($A \setminus B$)~---~разность документов (оператор \texttt{NOT});
	\item \textit{Adjacent}~---~документы, где термы стоят подряд (оператор \texttt{ADJ});
	\item \textit{Near}~---~документы, где термы стоят на расстоянии $\leq k$ (оператор \texttt{NEAR/$k$}).
\end{itemize}

\section{Skip lists}

Skip list~---~вероятностная структура данных поверх упорядоченного списка, позволяющая выполнять поиск за $O(log(n))$ вместо $O(n)$.
В данной реализации skip list строится детерминированно: на уровне $i$ размещаются элементы с шагом $2^{i+1}$, начиная с $\lceil\sqrt{n}\rceil$.
Итератор \texttt{PostingIterator} поддерживает курсоры на каждом уровне и метод \texttt{AdvanceTo(docID)}, который перемещает итератор к первому элементу $\geq docID$, используя skip-указатели для быстрого продвижения.

\section{Сжатие posting lists}

Если длина массива $\geq 32$, используется поблочный битпакинг (block size 32), где каждый блок из 32 дельт кодируется минимальным числом бит; иначе применяется varint-кодирование (\texttt{binary.PutUvarint}).
Битпакинг обеспечивает сжатие в 3--4 раза по сравнению с varint за счёт хранения фиксированного числа бит на элемент вместо переменного.

\section{Ранжирование BM25}

BM25~---~функция ранжирования, вычисляющая релевантность документа запросу:
\[
	\text{score}(D, Q) = \sum_{t \in Q} \text{IDF}(t) \cdot \frac{f(t, D) \cdot (k_1 + 1)}{f(t, D) + k_1 \cdot (1 - b + b \cdot |D| / \text{avgdl})}
\]

, где $f(t, D)$~---~частота терма $t$ в документе $D$, 
$|D|$~---~длина документа,
$\text{avgdl}$~---~средняя длина документа,
$k_1 = 1.5$ и $b = 0.75$~---~параметры.

$\text{IDF}(t) = \ln(1 + (N - n(t) + 0.5) / (n(t) + 0.5))$

, где $N$~---~число документов,
$n(t)$~---~число документов, содержащих терм $t$.

\chapter{Практическая часть}

\section{Архитектура реализации}

Проект реализован на Go и состоит из следующих пакетов:
\begin{enumerate}
	\item \texttt{internal/index}:
	\begin{itemize}
		\item \texttt{IndexBuilder}~---~построение индекса, добавление документов, токенизация, построение posting lists, дельта-кодирование и битпакинг;
		\item \texttt{InvertedIndex}~---~загрузка индекса через \texttt{mmap}, декодирование posting list;
		\item \texttt{PostingList}~---~постинг-лист со skip list и итератором \texttt{PostingIterator};
		\item операции \texttt{Intersect}, \texttt{Union}, \texttt{Difference}, \texttt{Adjacent}, \texttt{Near}.
	\end{itemize}
	\item \texttt{internal/query}:
	\begin{itemize}
		\item \texttt{Parser}~---~рекурсивный парсинг запроса с операторами \texttt{AND}, \texttt{OR}, \texttt{NOT}, \texttt{ADJ}, \texttt{NEAR/$k$} и скобками;
		\item \texttt{Evaluator}~---~обход AST и вызов соответствующих операций над posting lists.
	\end{itemize}
	\item \texttt{internal/ranking}~---~ранжирование BM25;
	\item \texttt{internal/compression}~---~битпакинг и varint-кодирование/декодирование posting lists;
	\item \texttt{internal/storage}~---~двоичный формат индекса (header + term entries + postings data) и хранилище документов (\texttt{DocStore});
	\item \texttt{internal/engine}~---~поисковый движок, объединяющий все необходимые компоненты;
	\item \texttt{cmd/tui}~---~TUI-интерфейс на \texttt{bubbletea}.
\end{enumerate}

\section{Формат хранения индекса}

Индекс хранится в двух файлах:
\begin{itemize}
	\item \texttt{*.idx}~---~сам индекс с заголовком (40 байт: magic, версия, число термов, число документов, число токенов, смещения), отсортированными term entries и сжатыми posting lists;
	\item \texttt{*.docs}~---~хранилище документов.
\end{itemize}

Обратный индекс загружается через \texttt{mmap} и читает только заголовок и term entries, а posting lists декодируются лениво при поиске.
Страницы \texttt{mmap} подгружаются операционной системой при первом обращении.
Для каждого терма posting list хранится в виде плоского массива: $[len_{docDeltas}][docDeltas][len_{posDeltas}][posDeltas][len_{tfs}][tfs][len_{skipList}][skipList]$, где каждый подмассив сжимается независимо (битпакинг или varint в зависимости от размера), а затем весь блок кодируется одним вызовом \texttt{Compress}.

\section{Поиск}

Обработка запроса состоит из трёх этапов:
\begin{enumerate}
	\item \textit{Парсинг}: строка запроса разбирается в AST;
	\item \textit{Вычисление}: обход AST, для каждого терма загружается posting list, операции выполняются в соответствии с приоритетом слева направо;
	\item \textit{Ранжирование}: для документов из результирующего posting list вычисляется BM25 score.
\end{enumerate}

\chapter{Исследовательская часть}

\section{Аппаратные характеристики}

Все замеры проводились на ноутбуке со следующей конфигурацией:
\begin{itemize}
	\item операционная система: \texttt{macOS} (arm64);
	\item процессор: \texttt{Apple M3 Pro}, 11 ядер;
	\item оперативная память: \texttt{18 ГБ};
	\item версия Go: \texttt{go1.25}.
\end{itemize}
Во время замеров не выполнялись ресурсоёмкие фоновые задачи, а ноутбук был подключён к сети электропитания.

\section{Методика замеров}

Для замеров использованы встроенные benchmark-тесты Go.
Измерялись четыре операции:
\begin{itemize}
	\item \texttt{build}~---~построение индекса;
	\item \texttt{warmup}~---~прогрев индекса;
	\item \texttt{search}~---~поиск по запросам различных типов (без ранжирования);
	\item \texttt{searchrank}~---~поиск с ранжированием BM25.
\end{itemize}

Для построения индекса использованы документы размером 500, 700, 1000, 2000, 3000, 4000, для поиска и warmup~---~документы размером 10000, 25000, 50000, 75000, 100000, 125000, 150000, 200000.
Разделение размеров обусловлено тем, что построение индекса при больших $N$ не успевает выполнить более одной итерации за таймаут.

Для каждого размера фиксировались:
\begin{itemize}
	\item среднее время на операцию (нс/оп);
	\item 95\% доверительный интервал (CI95);
	\item размер индекса в байтах (для build).
\end{itemize}

Поиск измеряется для следующих типов запросов: \texttt{Term}, \texttt{And}, \texttt{Or}, \texttt{Not}, \texttt{Adj}, \texttt{Near}, а также составных: \texttt{Complex\_AndOr}, \texttt{Complex\_AdjAnd}, \texttt{Complex\_AndNot}.
Для каждого типа генерируется 10000 случайных запросов.
В качестве термов используются термы из нижней половины распределения по частоте документа, что позволяет избежать аномально медленных запросов на высокочастотных термах.
Результатом является среднее время на один запрос.

\section{Результаты построения индекса}

Таблица~\ref{tab:build} содержит результаты замеров времени построения индекса и его размера.
\begin{table}[H]
	\centering
	\caption{Время построения индекса и размер}
	\label{tab:build}
	\begin{tabular}{|r|r|r|}
		\toprule
		$N$ & время (мс) & размер (МБ) \\
		\midrule
		500  & $281 \pm 14$  & 5  \\
		700  & $334 \pm 13$  & 7  \\
		1000 & $456 \pm 4$   & 9  \\
		2000 & $851 \pm 44$  & 16 \\
		3000 & $1206 \pm 27$ & 23 \\
		4000 & $1623 \pm 7$  & 29 \\
		\bottomrule
	\end{tabular}
\end{table}

График времени построения индекса представлен на рисунке~\ref{img:build_time}.
\includeimage
	{build_time}
	{f}
	{H}
	{\textwidth}
	{Время построения индекса в зависимости от числа документов}

Время построения растёт линейно.
При увеличении $N$ в 8 раз (от 500 до 4000) время увеличивается примерно в 6 раз (от 281 до 1623 мс).

\section{Результаты warmup}

Таблица~\ref{tab:warmup} содержит результаты замеров прогрева индекса.
\begin{table}[H]
	\centering
	\caption{Время warmup (мс)}
	\label{tab:warmup}
	\begin{tabular}{|r|r|}
		\toprule
		$N$ & время (мс) \\
		\midrule
		10000  & $129 \pm 16$  \\
		25000  & $366 \pm 23$  \\
		50000  & $644 \pm 13$  \\
		75000  & $744 \pm 60$  \\
		100000 & $823 \pm 8$   \\
		125000 & $852 \pm 32$  \\
		150000 & $1097 \pm 15$ \\
		200000 & $1337\pm 42$  \\
		\bottomrule
	\end{tabular}
\end{table}

График времени warmup представлен на рисунке~\ref{img:warmup_time}.
\includeimage
	{warmup_time}
	{f}
	{H}
	{\textwidth}
	{Время первого запроса после загрузки индекса (warmup)}

Время warmup значительно превышает время повторного поиска.
Это объясняется тем, что при первом обращении к posting lists операционная система подгружает страницы из диска в память.
С ростом $N$ время warmup растёт, что объясняется увеличением числа \texttt{mmap}-страниц, к которым обращается поисковый запрос.

\section{Результаты поиска в индексе}

Таблица~\ref{tab:search-basic} содержит результаты замеров времени базовых операций поиска.
\begin{table}[H]
	\centering
	\caption{Время базовых операций поиска (нс)}
	\label{tab:search-basic}
	\begin{tabular}{|r|r|r|r|r|r|r|}
		\toprule
		$N$ & Term & And & Or & Not & Adj & Near \\
		\midrule
		10000  & $2771 \pm 3$    & $5821 \pm 5$    & $6073 \pm 5$     & $6315 \pm 5$   & $6004 \pm 5$   & $6138 \pm 5$  \\
		25000  & $2876 \pm 3$    & $5943 \pm 5$    & $6413 \pm 7$     & $6398 \pm 5$   & $6105 \pm 5$   & $6276 \pm 6$  \\
		50000  & $2916 \pm 3$    & $6063 \pm 5$    & $7135 \pm 58$    & $6461 \pm 6$   & $6213 \pm 5$   & $6324 \pm 6$  \\
		75000  & $2926 \pm 3$    & $6179 \pm 6$    & $10385 \pm 137$  & $6548 \pm 5$   & $6262 \pm 6$   & $6394 \pm 6$  \\
		100000 & $2952 \pm 3$    & $6514 \pm 7$    & $11700 \pm 398$  & $6660 \pm 6$   & $6329 \pm 6$   & $6426 \pm 6$  \\
		125000 & $3064 \pm 3$    & $6512 \pm 7$    & $12832 \pm 884$  & $6720 \pm 6$   & $6383 \pm 6$   & $6485 \pm 6$  \\
		150000 & $3055 \pm 3$    & $7041 \pm 9$    & $16720 \pm 383$  & $6911 \pm 6$   & $6504 \pm 6$   & $6702 \pm 6$  \\
		200000 & $3143 \pm 3$    & $7134 \pm 9$    & $26040 \pm 802$  & $6856 \pm 6$   & $6602 \pm 6$   & $6767 \pm 6$  \\
		\bottomrule
	\end{tabular}
\end{table}

Таблица~\ref{tab:search-complex} содержит результаты замеров времени составных операций поиска.
\begin{table}[H]
	\centering
	\caption{Время составных операций поиска (нс)}
	\label{tab:search-complex}
	\begin{tabular}{|r|r|r|r|}
		\toprule
		$N$ & OrAnd & AdjAnd & AndNot \\
		\midrule
		10000  & $12060 \pm 11$ & $9009 \pm 8$   & $12330 \pm 10$ \\
		25000  & $12335 \pm 25$ & $9220 \pm 8$   & $12536 \pm 11$ \\
		50000  & $12264 \pm 10$ & $9277 \pm 8$   & $12725 \pm 11$ \\
		75000  & $12837 \pm 59$ & $9406 \pm 8$   & $12720 \pm 11$ \\
		100000 & $12765 \pm 51$ & $9485 \pm 8$   & $12799 \pm 11$ \\
		125000 & $12923 \pm 53$ & $9592 \pm 8$   & $12866 \pm 11$ \\
		150000 & $13221 \pm 77$ & $9635 \pm 8$   & $12956 \pm 11$ \\
		200000 & $13289 \pm 63$ & $9999 \pm 9$   & $13105 \pm 12$ \\
		\bottomrule
	\end{tabular}
\end{table}

Графики времени поиска представлены на рисунках~\ref{img:search_latency} и~\ref{img:search_complex}.
\includeimage
	{search_latency}
	{f}
	{H}
	{\textwidth}
	{Время базовых операций поиска}

\includeimage
	{search_complex}
	{f}
	{H}
	{\textwidth}
	{Время составных операций поиска}

\section{Результаты поиска с ранжированием}

Таблица~\ref{tab:searchrank-basic} содержит результаты замеров времени базовых операций поиска с ранжированием BM25.
\begin{table}[H]
	\centering
	\caption{Время базовых операций поиска с ранжированием (нс)}
	\label{tab:searchrank-basic}
	\begin{tabular}{|r|r|r|r|r|r|r|}
		\toprule
		$N$ & Term & And & Or & Not & Adj & Near \\
		\midrule
		10000  & $2960 \pm 3$    & $5879 \pm 5$    & $6398 \pm 6$     & $6512 \pm 5$   & $6079 \pm 5$   & $6203 \pm 5$  \\
		25000  & $3021 \pm 3$    & $6069 \pm 5$    & $6575 \pm 6$     & $6662 \pm 6$   & $6181 \pm 5$   & $6358 \pm 6$  \\
		50000  & $3082 \pm 3$    & $6185 \pm 5$    & $7165 \pm 22$    & $6740 \pm 6$   & $6256 \pm 6$   & $6433 \pm 6$  \\
		75000  & $3106 \pm 3$    & $6257 \pm 6$    & $10381 \pm 485$  & $6821 \pm 6$   & $6306 \pm 6$   & $6517 \pm 6$  \\
		100000 & $3231 \pm 3$    & $6713 \pm 9$    & $13234 \pm 957$  & $6981 \pm 6$   & $6360 \pm 6$   & $6543 \pm 6$  \\
		125000 & $3259 \pm 3$    & $6609 \pm 7$    & $14449 \pm 1091$ & $6986 \pm 6$   & $6409 \pm 6$   & $6626 \pm 6$  \\
		150000 & $3535 \pm 4$    & $7560 \pm 11$   & $20583 \pm 1275$ & $7415 \pm 7$   & $6895 \pm 7$   & $7034 \pm 7$  \\
		200000 & $3674 \pm 4$    & $8218 \pm 38$   & $28927 \pm 1099$ & $7473 \pm 7$   & $6959 \pm 7$   & $7153 \pm 7$  \\
		\bottomrule
	\end{tabular}
\end{table}

Таблица~\ref{tab:searchrank-complex} содержит результаты замеров времени составных операций поиска с ранжированием.
\begin{table}[H]
	\centering
	\caption{Время составных операций поиска с ранжированием (нс)}
	\label{tab:searchrank-complex}
	\begin{tabular}{|r|r|r|r|}
		\toprule
		$N$ & OrAnd & AdjAnd & AndNot \\
		\midrule
		10000  & $12102 \pm 10$ & $9181 \pm 8$   & $12314 \pm 10$ \\
		25000  & $12333 \pm 10$ & $9396 \pm 8$   & $12568 \pm 11$ \\
		50000  & $12511 \pm 11$ & $9506 \pm 8$   & $12688 \pm 11$ \\
		75000  & $12844 \pm 48$ & $9686 \pm 8$   & $12768 \pm 11$ \\
		100000 & $12945 \pm 52$ & $9735 \pm 8$   & $12812 \pm 11$ \\
		125000 & $13074 \pm 51$ & $9754 \pm 8$   & $12931 \pm 11$ \\
		150000 & $13725 \pm 57$ & $10238 \pm 10$ & $13455 \pm 14$ \\
		200000 & $14464 \pm 83$ & $10507 \pm 11$ & $13675 \pm 14$ \\
		\bottomrule
	\end{tabular}
\end{table}

Графики времени поиска с ранжированием BM25 представлены на рисунках~\ref{img:search_rank_latency} и~\ref{img:search_rank_complex}.
\includeimage
	{search_rank_latency}
	{f}
	{H}
	{\textwidth}
	{Время базовых операций поиска с ранжированием}

\includeimage
	{search_rank_complex}
	{f}
	{H}
	{\textwidth}
	{Время составных операций поиска с ранжированием}

Время поиска растёт с увеличением $N$ для всех операторов, что объясняется ростом длины posting lists и, соответственно, увеличением объёма вычислений при пересечении/объединении.

Оператор \texttt{Term} является самым быстрым ($\approx$3 мкс), поскольку требует лишь загрузки одного posting list.
Среди бинарных операторов \texttt{And} и \texttt{Not}~---~самые быстрые ($\approx$7 мкс при $N = 200000$), а \texttt{Or}~---~самый медленный ($\approx$26 мкс).
Операторы \texttt{Adj} и \texttt{Near} медленнее \texttt{And}.

Составные запросы показывают время, сопоставимое с суммой времён базовых операторов, что подтверждает корректность поэтапного вычисления.
\texttt{Complex\_OrAnd}~---~самый медленный составной запрос ($\approx$13 мкс при $N = 200000$), что объясняется наличием операции \texttt{Or}.

\section{Размер индекса}

График размера индекса в зависимости от числа документов представлен на рисунке~\ref{img:index_size}.
\includeimage
	{index_size}
	{f}
	{H}
	{\textwidth}
	{Размер индекса в зависимости от числа документов}

Размер индекса растёт линейно с числом документов.
Это объясняется тем, что основную долю занимают posting lists и данные документов, объём которых пропорционален числу токенов.

\section{Профилирование}

Для профилирования используется отдельная программа \texttt{cmd/profile/main.go}, которая выполняет построение индекса или поисковые запросы с включённым \texttt{pprof.StartCPUProfile}.

На рисунках~\ref{img:build_profile} и~\ref{img:search_profile} представлены графы вызовов CPU-профиля для построения индекса и поиска соответственно.
\includeimage
	{build_profile}
	{f}
	{H}
	{\textwidth}
	{CPU-профиль построения индекса}

\includeimage
	{search_profile}
	{f}
	{H}
	{\textwidth}
	{CPU-профиль поиска}

В профиле построения индекса основное время тратится на:
\begin{itemize}
	\item \texttt{IndexBuilder.AddDocument}~---~токенизацию и обновление posting lists;
	\item \texttt{IndexBuilder.Save}~---~сортировку термов, битпакинг, построение skip lists и запись на диск.
\end{itemize}

В профиле поиска основное время приходится на:
\begin{itemize}
	\item \texttt{InvertedIndex.GetPostings}~---~бинарный поиск терма и декодирование posting list;
	\item \texttt{Intersect}/\texttt{Union}/\texttt{Adjacent}/\texttt{Near}~---~операции над posting lists;
	\item \texttt{Ranker.RankResults}~---~вычисление BM25-оценок с использованием итераторов.
\end{itemize}

\chapter{Вывод}

Реализован обратный индекс с поддержкой поисковых операторов \texttt{AND}, \texttt{OR}, \texttt{NOT}, \texttt{ADJ}, \texttt{NEAR/$k$}, ранжированием BM25 и TUI-интерфейсом.
Индекс хранится на диске в бинарном формате с битпакингом и дельта-кодированием posting lists и загружается через \texttt{mmap} с ленивым декодированием.

Проведены замеры производительности на наборах данных от 500 до 200000 документов.
Можно сделать следующие выводы:
\begin{itemize}
	\item время построения индекса растёт линейно;
	\item время warmup значительно превышает время повторного поиска;
	\item оператор \texttt{Term}~---~самый быстрый; \texttt{Or}~---~самый медленный;
	\item размер индекса растёт линейно.
\end{itemize}

\end{document}
